\textbf{Binary detection.} On the Experiment~1 test set the single tree ensembles are the strongest binary detectors (Table~\ref{tab:binary}): XGBoost reaches \EoneXgbRecall\% recall at \EoneXgbFpr\% FPR (\EoneXgbAlerts{} alerts per 1000 benign flows), and Random Forest \EoneRfRecall\% at \EoneRfFpr\%. The tuned SVM gate reaches its designed recall (\EoneGateRecall\%) but at a higher false-positive rate (\EoneGateFpr\%, ROC-AUC \EoneGateAuc), and Logistic Regression is the weakest (\EoneLrFpr\% FPR). The fused system inherits the gate's recall and adds the autoencoder's anomaly alerts, giving \EoneFusionRecall\% recall at \EoneFusionFpr\% FPR. The multi-stage design therefore does not beat a well-configured tree ensemble on raw binary accuracy in this closed-set setting. Its value is in attribution, independent anomaly evidence and explainability, and its binary performance is limited by the gate.

\textbf{Family attribution.} All four models attribute families of true attack flows with more than 99.5\% accuracy (Table~\ref{tab:multiclass}). On macro-F1, which weights the six-flow Infiltration class equally with DoS, the Random Forest base learner (\EoneRfMacroFone) scores higher than the stacked ensemble (\EoneStackMacroFone), XGBoost (\EoneXgbMacroFone) and Logistic Regression (\EoneLrMacroFone). Stacking therefore did not improve over its best base model here. Its meta-learner is trained with balanced class weights on out-of-fold probabilities, and with only a few dozen Infiltration and Heartbleed rows the meta-learner's estimates for these classes are noisy.

\textbf{Generalisation to an unseen capture.} In Experiment~2 the test file is the Friday afternoon DDoS capture, and DDoS never appears in training. The trained models behave very differently there (Table~\ref{tab:experiments}). The tree ensembles keep their precision but detect only \EtwoRfRecall\% (Random Forest) and \EtwoXgbRecall\% (XGBoost) of the unseen DDoS flows. The SVM gate's threshold was calibrated on the only validation file, which contains web attacks. On the new day it flags almost all traffic: \EtwoGateRecall\% of attacks but also \EtwoGateFpr\% of benign flows. The autoencoder's benign false-positive rate rises from \EoneAeFpr\% to \EtwoAeFpr\%, which shows that benign traffic itself shifts between capture days. When the family classifier is given true DDoS flows, it labels them as the closest known families: about 84\% as DoS and 16\% as PortScan (from the Experiment~2 confusion matrix).

